% MAEG4060 — Stage 1 Proposal (recreated from PDF text)
\documentclass[11pt,a4paper]{article}

\usepackage[margin=2.5cm]{geometry}
\usepackage[T1]{fontenc}
\usepackage{lmodern}
\usepackage{microtype}
\usepackage{graphicx}
\usepackage{booktabs}
\usepackage{amsmath}
\usepackage{url}
\usepackage[hidelinks,breaklinks]{hyperref}

\makeatletter
\g@addto@macro\UrlBreaks{\do\/\do-\do\.\do\=}
\makeatother

\newcommand{\figureplaceholder}[2]{%
  \begin{figure}[htbp]
    \centering
    \fbox{%
      \parbox[c][5cm][c]{0.88\linewidth}{%
        \centering\small\textit{#1}\\[0.5em]#2%
      }%
    }
  \end{figure}%
}

\title{%
  MAEG4060 Virtual Reality Systems and Applications\\
  \large Course Project -- Stage 1 Proposal (10\%)}
\author{%
  \parbox{0.92\linewidth}{\centering
  \textbf{Project Title:} Interactive Physical Simulation of Off-Road Vehicle Dynamics in a Mountainous Terrain\\[1em]
  \textbf{Group Members:}\\[0.35em]
  Guo Tsun Hei (1155186272): 3D Modeling, Texturing, Rigging, and Animation\\
  Yim Fu Chong (1155176974): Programming, Physics Engine, and System Integration}}
\date{Submission deadline: 23 Jan 2026}

\begin{document}
\maketitle

\section{Project Theme and Objectives}

The objective of this project is to design and develop a high-fidelity, interactive 3D ``Physics Sandbox'' that simulates the mechanical relationship between a rigid-body vehicle and a complex, non-planar terrain.

In strict adherence to the course requirements, this application is not a video game; it contains no missions or scoring systems. Instead, it serves as a technical demonstration of:

\begin{enumerate}
  \item \textbf{Kinematic Terrain-Following:} Real-time calculation of vehicle orientation (Pitch/Roll) based on heightmap gradients.
  \item \textbf{Momentum \& Friction Simulation:} Accurate modelling of acceleration, braking, and rolling resistance.
  \item \textbf{Dynamic Environmental Interaction:} A living environment featuring hierarchical animations and a time-of-day system.
\end{enumerate}

The application will be developed in Visual C++ using the DirectX 9.0c Fixed Function Pipeline.

\section{Detailed Model Specifications (20\%)}

\textit{Submission deadline: 27 Feb 2026}

All assets will be exported in the \texttt{.X} file format to ensure compatibility with DirectX frame hierarchies. The project utilizes a Left-Handed (LH) Coordinate System (consistent with standard \texttt{D3DXMatrixLookAtLH} functions).

\subsection{Landscape/Terrain}

\begin{itemize}
  \item \textbf{Geometry:} Generated via a $1024\times 1024$ Heightmap, providing a high-resolution mesh for smooth physics calculations.
  \item \textbf{Surface Shading:} Implementation of Texture Splatting. A custom algorithm will blend ``Lush Grass'' textures on flat plains with ``Cragged Rock'' textures on slopes exceeding a 30-degree incline.
  \item \textbf{Visual Reference -- FIGURE 1:} Sketch of Mountain/Valley Layout
\end{itemize}

\figureplaceholder{FIGURE 1}{Sketch illustrating the elevation variance, including the central valley for vehicle testing and the steep cliffs for texture blending demonstrations.}

\noindent\textbf{Source:} \url{https://sketchfab.com/3d-models/half-dome-terrain-landscape-86ef075dc7d9414583ac2180e1e651a3}

\subsection{Controllable Vehicle (Jeep/ATV)}

\begin{itemize}
  \item \textbf{Hierarchy:} The model features a parent body mesh with four independent wheel nodes. This allows for procedural rotation and steering transformations relative to the chassis.
  \item \textbf{Efficiency:} Poly-count is optimized at approximately 3{,}000 triangles to ensure high-frequency collision detection without CPU bottlenecking.
  \item \textbf{Visual Reference -- FIGURE 2:} Sketch of Vehicle Model
\end{itemize}

\figureplaceholder{FIGURE 2}{Technical sketch showing the hierarchical structure (Wheel\_FL, Wheel\_FR, etc.) and the low-poly chassis design.}

\noindent\textbf{Source:} \url{https://sketchfab.com/3d-models/atv-from-game-pure-f2e454611ebd4151b9d7aa31815c35e4}

\subsection{Animated Environmental Elements (20\%)}

\begin{itemize}
  \item \textbf{Maya Keyframed Animation:} To satisfy the group project requirement, a Windmill prop will be modeled and animated within Maya. The sail rotation will be exported as an AnimationSet in the \texttt{.X} file and played back in the engine using \texttt{ID3DXAnimationController}.
  \item \textbf{Procedural Animation:} Vehicle wheels will rotate dynamically based on the calculated linear velocity ($v = \omega r$) to ensure visual fidelity during acceleration and braking.
  \item \textbf{Visual Reference -- FIGURE 3:} Sketch of Windmill/Props
\end{itemize}

\begin{figure}[htbp]
  \centering
  \fbox{%
    \parbox[c][5cm][c]{0.88\linewidth}{%
      \centering\small\textit{FIGURE 3}\\[0.5em]Sketch of Windmill/Props (reference placement)%
    }%
  }
  \caption{Sketch of the windmill sub-mesh components (target for Maya animation) and environmental props like crates and boulders.}
\end{figure}

\noindent\textbf{Source:} \url{https://sketchfab.com/3d-models/rock-boulder-game-ready-asset-0079cb0e717c4aeebe6bca21ec21b0d9}

\noindent\textbf{Source:} \url{https://sketchfab.com/3d-models/windmill-fe13862b9e3c4dc887d10e297d86f17f}

\noindent\textbf{Source:} \url{https://sketchfab.com/3d-models/wooden-crate-3c0259d17f3b4306890cdc809b545670}

\section{Implementation of Effects and Physics}

\subsection{Movement \& Control (5\%)}

\begin{itemize}
  \item \textbf{Input Handling:} Integration of DirectInput for responsive WASD keyboard control and XInput for analog gamepad support (trigger-based throttle).
  \item \textbf{Dynamics:} Implementation of inertia and rolling friction constants to simulate realistic deceleration when the throttle is released.
\end{itemize}

\subsection{Terrain Navigation (5\%)}

\begin{itemize}
  \item \textbf{Bilinear Interpolation:} To prevent ``snapping,'' the vehicle's Y-height will be calculated by interpolating between the four nearest heightmap grid points.
  \item \textbf{Surface Normal Alignment:} The vehicle's ``Up'' vector will be synchronized with the terrain's surface normal (calculated via the cross-product of the current triangle vertices).
\end{itemize}

\subsection{Collision Detection \& Response (10\% + 10\% Bonus)}

\begin{itemize}
  \item \textbf{Static Collision:} AABB (Axis-Aligned Bounding Boxes) will be used for rectangular props (crates), while Bounding Spheres will be used for irregular objects (boulders).
  \item \textbf{Bonus Objective:} Dynamic-to-Dynamic Collision. We will implement a momentum transfer system where the vehicle can displace ``loose'' boulders, causing them to roll down inclines and interact with other world objects.
\end{itemize}

\subsection{Physics-based Simulation \& Natural Phenomena (10\%)}

\begin{itemize}
  \item \textbf{Gravity:} A constant downward force vector ($g = 9.81$ units/s$^2$) applied to all non-static entities.
  \item \textbf{Day/Night Cycle:}
    \begin{itemize}
      \item \textbf{Solar Rotation:} The \texttt{D3DLIGHT9} directional light source will orbit the scene using $\sin(\text{time})$ and $\cos(\text{time})$ functions to update its direction vector every frame.
      \item \textbf{Atmospheric Scattering:} The background clear color and light diffuse/ambient properties will transition from bright amber (Day) to deep indigo (Night).
    \end{itemize}
\end{itemize}

\subsection{Visual Enhancements \& User Guide (5\%)}

\begin{itemize}
  \item \textbf{On-Screen Display (OSD):} Utilizing \texttt{ID3DXFont} to provide a real-time ``Debug HUD'' showing current velocity, world coordinates $(X, Y, Z)$, and the current state of the Day/Night cycle.
\end{itemize}

\section{Work Division}

\begin{table}[htbp]
  \centering
  \begin{tabular}{@{}p{0.18\linewidth} p{0.18\linewidth} p{0.22\linewidth} p{0.36\linewidth}@{}}
    \toprule
    \textbf{Task} & \textbf{Category} & \textbf{Primary Member} & \textbf{Technical Detail} \\
    \midrule
    Stage 2: & Assets & Guo Tsun Hei & Maya modeling, UV Unwrapping, \texttt{.X} export, hierarchical rigging, and creating keyframe animations in Maya. \\
    Stage 3: & Coding & Yim Fu Chong & DX9 Device initialization, Heightmap parser, Physics engine, Collision logic, and Day/Night system. \\
    & Integration & Both & Debugging mesh loading errors and tuning physics constants for ``feel.'' \\
    \bottomrule
  \end{tabular}
\end{table}

\section{Final Report Planning (15\%)}

\textit{Submission deadline: 10 April 2026}

The final submission will include:

\begin{enumerate}
  \item \textbf{The Executable \& Source:} A completed, compiled Visual C++ application.
  \item \textbf{User Guide:} Detailed instructions for controls and hardware requirements.
  \item \textbf{Technical Description:} A brief description of the mathematical techniques used for the physics engine, collision detection, and terrain navigation.
\end{enumerate}

\section{References \& Asset Credits}

\begingroup
\sloppy
\setlength{\emergencystretch}{4em}

\begin{enumerate}
  \item \textbf{Half Dome Terrain Landscape}
    \begin{itemize}\setlength\itemsep{2pt}\setlength\parskip{0pt}
      \item Author: Studio Lab
      \item License: Creative Commons Attribution (CC BY 4.0)
      \item Source: Sketchfab
      \item URL: \url{https://sketchfab.com/3d-models/half-dome-terrain-landscape-86ef075dc7d9414583ac2180e1e651a3}
      \item Usage: Reference for high-resolution terrain geometry and heightmap generation.
    \end{itemize}

  \item \textbf{Controllable Vehicle (ATV)}
    \begin{itemize}\setlength\itemsep{2pt}\setlength\parskip{0pt}
      \item Title: ATV from game ``PURE''
      \item Author: barking\_dogo
      \item License: Creative Commons Attribution (CC BY 4.0)
      \item Source: Sketchfab
      \item URL: \url{https://sketchfab.com/3d-models/atv-from-game-pure-f2e454611ebd4151b9d7aa31815c35e4}
      \item Usage: Main controllable vehicle model (Wheel nodes separated for procedural animation).
    \end{itemize}

  \item \textbf{Rock Boulder (Game Ready Asset)}
    \begin{itemize}\setlength\itemsep{2pt}\setlength\parskip{0pt}
      \item Author: Aparicio Silva 3D
      \item License: Creative Commons Attribution (CC BY 4.0)
      \item Source: Sketchfab
      \item URL: \url{https://sketchfab.com/3d-models/rock-boulder-game-ready-asset-0079cb0e717c4aeebe6bca21ec21b0d9}
      \item Usage: Used for dynamic physics simulation (Bounding Sphere collision).
    \end{itemize}

  \item \textbf{Windmill}
    \begin{itemize}\setlength\itemsep{2pt}\setlength\parskip{0pt}
      \item Author: KIFIR
      \item License: Creative Commons Attribution (CC BY 4.0)
      \item Source: Sketchfab
      \item URL: \url{https://sketchfab.com/3d-models/windmill-fe13862b9e3c4dc887d10e297d86f17f}
      \item Usage: Used for keyframed animation and environmental aesthetics.
    \end{itemize}

  \item \textbf{Wooden Crate}
    \begin{itemize}\setlength\itemsep{2pt}\setlength\parskip{0pt}
      \item Author: Domingos Studios
      \item License: Creative Commons Attribution (CC BY 4.0)
      \item Source: Sketchfab
      \item URL: \url{https://sketchfab.com/3d-models/wooden-crate-3c0259d17f3b4306890cdc809b545670}
      \item Usage: Used for static collision detection (AABB).
    \end{itemize}
\end{enumerate}

\endgroup

\appendix
\section{Project Execution \& Technical Roadmap}

\subsection{Core Project Concept}

This project is explicitly defined as a ``3D Physics Sandbox'' rather than a traditional video game. It deviates from mission-based or score-based gameplay to focus entirely on technical simulation. The core objective is to demonstrate a realistic mechanical interaction between a rigid-body off-road vehicle and complex, non-planar mountain terrain, governed by fundamental physics laws (gravity, friction, and inertia).

\subsection{Role Responsibilities \& Workflow}

To ensure efficiency and strict adherence to the Group Project requirements, the workload is divided as follows:

\paragraph{Guo Tsun Hei (Assets, Animation \& Pipeline)}
\begin{itemize}
  \item \textbf{Primary Responsibility:} Creating all visual assets and ensuring they are technically compatible with the DirectX pipeline.
  \item \textbf{Key Deliverables:}
    \begin{itemize}
      \item \textbf{Maya Animation (Critical Requirement):} Creation of the ``Windmill'' prop with keyframed rotation animations. This serves as the mandatory pre-baked animation asset.
      \item \textbf{Vehicle Hierarchy:} Modeling the Jeep with four independent wheel nodes (children of the chassis) to allow for procedural rotation control in the engine.
      \item \textbf{Pipeline Management:} Exporting all scenes (Terrain, Vehicle, Props) into the \texttt{.X} format for DirectX 9 integration.
    \end{itemize}
\end{itemize}

\paragraph{Yim Fu Chong (Engine, Physics \& Integration)}
\begin{itemize}
  \item \textbf{Primary Responsibility:} Bringing static assets to life through code and physics calculations.
  \item \textbf{Key Deliverables:}
    \begin{itemize}
      \item \textbf{Kinematic Terrain Following:} Implementing algorithms to calculate heightmap gradients, ensuring the vehicle ``sticks'' to the ground while adjusting Pitch/Roll orientation dynamically.
      \item \textbf{Physics Simulation:} Coding the forces for Gravity ($g = 9.81$), friction, and momentum conservation (preventing instant stops).
      \item \textbf{Procedural Animation:} Implementing code-driven wheel rotation ($v = \omega r$) that corresponds to vehicle velocity.
      \item \textbf{Environmental Logic:} Implementing the Day/Night cycle using trigonometric light orbit functions.
    \end{itemize}
\end{itemize}

\subsection{Technical Highlights \& Bonus Objectives}

The project aims to demonstrate proficiency in three specific areas:

\begin{itemize}
  \item \textbf{Advanced Texturing:} Implementation of Texture Splatting, where the terrain shader automatically switches between ``Grass'' and ``Rock'' textures based on a slope threshold ($>30$ degrees).
  \item \textbf{Hybrid Collision Strategy:} Demonstrating knowledge of multiple algorithms by using AABB (Axis-Aligned Bounding Boxes) for crates and Bounding Spheres for boulders.
  \item \textbf{Bonus Objective (Dynamic-to-Dynamic Interaction):} A momentum transfer system where the vehicle can collide with ``loose'' boulders, causing them to physically roll down the terrain, rather than acting as static walls.
\end{itemize}

\subsection{Project Status (as of Jan 20, 2026)}

The Stage 1 Proposal is complete. The team has established the technical scope and asset requirements. The workflow is finalized: Maya will be used for high-fidelity modeling and keyframed animation, while Visual C++ / DirectX 9 will handle the real-time physics simulation and procedural effects. The team is prepared for full development immediately following this submission.

\end{document}
